\subsection{RQ3: Is the gap an artifact of tuning, budget, shaping or the observation?}
\label{sec:rq3}
% Every number is generated; a guarded block prints "pending" until its
% experiment has finished, so the paper compiles at every stage.
\newcommand{\ifmacro}[3]{\ifcsname #1\endcsname #2\else #3\fi}

\paragraph{PPO tuning (E3).} Eight one-factor variations of PPO's
configuration were trained on root 42 and compared with the default on
validation seeds only (Appendix~\ref{app:hyper}).
\ifmacro{TuneEThreeRuns}{Of the \TuneEThreeRuns{} that have finished, the
best by validation return is \texttt{\TuneBest} (\TuneBestVal, against
\TuneDefaultVal{} for the default and \TuneBanditValBest{} for the bandit on
the same root); Table~\ref{tab:tuning} lists all of them with their test
returns, which played no role in selection.}{(pending)}

\paragraph{Budget (E5).} \ifmacro{RQThreeBudgetPPOvsBandit}{With three times
the budget (1.2M transitions, root 42), PPO reaches \RQThreeBudgetPPOvsBanditA{}
against \RQThreeBudgetPPOvsBanditB{} for the bandit at 400k (difference
\RQThreeBudgetPPOvsBandit).\ifmacro{RQThreeBudgetBandit}{ The bandit itself
changes by \RQThreeBudgetBandit{} with the larger budget.}{}}{(pending)}

\paragraph{Shaping (E6).} The bandit is trained on the shaped reward, and
potential-based shaping is policy-invariant only for the discount it was
designed for; for a myopic learner it adds one step of lookahead on the
potential. \ifmacro{RQThreeNoShaping}{Training the bandit with the shaping
term removed and evaluating it on the primary reward changes its return by
\RQThreeNoShaping{} (\RQThreeNoShapingRootsPositive{} of
\RQThreeNoShapingRoots{} roots higher), so the bandit's advantage does not rest
on the shaping term.}{(pending)}

\paragraph{Time of day (E7).} The open-loop diagnostic attributes the
non-myopic value to anticipation of exogenous change, and the observation has
no clock. If that is the missing ingredient for a sequential learner, adding
the time of day should help a $\gamma=0.9$ Q-learner more than the bandit
(H9, exploratory). \ifmacro{RQThreeTimeQ}{Adding the clock changes the
bandit's return by \RQThreeTimeBandit{} (\RQThreeTimeBanditRootsPositive{} of
\RQThreeTimeBanditRoots{} roots higher) and the Q-learner's by
\RQThreeTimeQ{} (\RQThreeTimeQRootsPositive{} of \RQThreeTimeQRoots{}
roots higher).}{(pending)}

\begin{table}[t]
\centering\small
\caption{Control experiments: per-root mean test returns of the compared
policies on their common roots, and the paired difference (mean; per root).}
\label{tab:rq3}
\begin{adjustbox}{max width=\columnwidth}
\begin{tabular}{@{}lrrrr@{}}
\toprule
Comparison & Roots & A & B & A $-$ B \\
\midrule
\IfFileExists{tables/rq3_controls.tex}{\tabinput{tables/rq3_controls}}{\multicolumn{5}{c}{pending}\\}
\bottomrule
\end{tabular}
\end{adjustbox}
\end{table}
